\medskip\noindent
\textbf{A source decomposition that treats every nine-cell at once.}
We now use the target conclusions to extend the ordinary map,
keeping the choices made through dimension eight explicit.
Put $Q_i=K_i/K_4$ as in D.2.18--D.2.19, and work with the
two-local quotient cell models there. The five eight-cell
attachments in \textup{(D.2.19w)} give, by cell slides,
\[
 Q_8\simeq D'\vee M(\ZZ/2,7)_F\vee M(\ZZ/4,7)_J
                  \vee C_1^8\vee S^8_s,
 \tag{D.2.20s1}
\]
where
\[
\begin{split}
 D&=\operatorname{cofib}\left(
  \begin{pmatrix}4&\eta_5\\0&2\end{pmatrix}:
  S^5\vee S^6\longrightarrow S^5\vee S^6\right),\\
 D'&=D\cup_{\tau_d+\tau_e}e^8_H,\qquad
 C_1^8=\operatorname{cofib}\bigl((2,\eta_6):S^6\vee S^7\to S^6\bigr).
\end{split}
 \tag{D.2.20s2}
\]
Indeed slide the cell $G_8$ over $H_8$, leaving the
degree-two attachment on $F$; slide $B_0^2$ over $A_0C$,
leaving the degree-four attachment on $J$; and slide
$Q^5A_0$ over $A_0C$, leaving the free eight-sphere.
The integral orientation of the last sphere is
$s=Q^5A_0-A_0C$. The cell $A_0C$ supplies the top
cell of $C_1^8$, while $H_8$ supplies the top cell of $D'$.
These are actual comparisons for the mapping problem:
the quotients are four-connected, their dimensions are
at most eight, and Freudenthal realizes the stable cell
slides unstably. Maps into the two-local space $T$ may
then be pulled back to the original finite skeleta.

In D.2.19 the correction is zero on the source six-sphere
$z$ and on the free seven-spheres $F,J$. Its nullhomotopies
on the eight-cells may therefore be chosen so that the
correction is constant on all the factors of
\textup{(D.2.20s1)} except $D'$. This makes its restrictions
to the two Moore factors, $C_1^8$, and the last sphere
literally zero. On $D'$ it retains the Moore extension
$\mu$, the bottom class $\overline\alpha$, and the
six-sphere value $\rho$. Its extension across the
$H$-cell can be any of the allowed nullhomotopies;
we will not assume that its effect on the next
homotopy group is an arbitrary free choice.

\medskip\noindent
\textbf{Every class in the remaining attaching group has zero
mod-two Hurewicz image.}
In $D$, the two classes $\tau_d,\tau_e$ are independent
order-two elements, by the complete calculation
\textup{(D.2.19p)}. Thus $\tau_d+\tau_e$ is nonzero
and has order two. The cofiber sequence
\[
 S^7\xrightarrow{\tau_d+\tau_e}D\longrightarrow D'
             \xrightarrow{p_H}S^8
 \tag{D.2.20s3}
\]
gives
\[
 \operatorname{im}\bigl((p_H)_*:\pi_8D'\to\pi_8S^8\bigr)
                  =2\ZZ_{(2)}.
 \tag{D.2.20s4}
\]
This is first computed in the stable exact sequence;
four-connectivity identifies all the relevant groups
with their ordinary counterparts. The complex $D$
has dimension seven and $H_7(D;\ZZ_{(2)})=0$:
its top cellular differential has nonzero degree two
on the free six-cell. The homology sequence of
\textup{(D.2.20s3)} consequently identifies
$H_8(D';\ZZ_{(2)})$ with $\ZZ_{(2)}$, by the
pinch onto the $H$-cell. It follows from
\textup{(D.2.20s4)} that every actual sphere map
$S^8\to D'$ has even degree on this sole integral
eight-dimensional homology generator. Therefore
\[
 h_{\FF_2}:\pi_8D'\longrightarrow H_8(D';\FF_2)
                  \quad\text{is zero}.
 \tag{D.2.20s5}
\]
There is no additional Tor summand in this last
homology group, since $H_7D'=0$. This argument
includes the entire torsion subgroup of $\pi_8D'$;
it does not infer its vanishing from incomplete
primary-square detectors.

\medskip\noindent
\textbf{The correction kills all actual nine-cell attachments.}
Let $f:D'\to T$ be the correction, with any of its
allowed fillings of the $H$-cell. It lifts to
$F=\fib(T\to B)$, since
$\Map_*(D',K(\FF_2,2))$ is weakly contractible.
Naturality and \textup{(D.2.20s5)} make the image
of every element of $\pi_8D'$ have zero mod-two
Hurewicz in $F$. The target result proved above
makes the actual Hurewicz homomorphism on $\pi_8F$
injective. Thus
\[
 f_*:\pi_8D'\longrightarrow\pi_8T
                       \quad\text{is zero}.
 \tag{D.2.20s6}
\]
The correction is constant on the other wedge
factors in \textup{(D.2.20s1)}. There are no cross
Whitehead terms in dimension eight by their
connectivities. Consequently its homomorphism
on all of $\pi_8Q_8$ is zero. Every nine-cell
attaching map has zero image, regardless of any
component not measured by primary operations.
Choose nullhomotopies on these cells to extend
the correction to $Q_9$.

Finally, the canonical summand already constructed
in \textup{(D.2.19y)--(D.2.19z)} extends through
source relative dimension nine. We keep that same
map and add the pullback of the correction along
$K_9\to Q_9$. This gives an actual ordinary map
$G_9:K_9\to T$ extending the allowed $G_8$ and
retaining every six-dimensional value in D.2.17.
The canonical summand has only been supplied
through dimension nine, and no extension over
the tenth or later cells follows from this proof.
